% !PRINTED 702
\weberpage{702}
%<<CONT
ein Differential dritter Gattung. Damit ist der Satz 4. bewiesen, sogar mit der Verschärfung, daß die darin auftretenden Differentiale dritter Gattung einen beliebigen festen Punkt $\mathfrak{P}_0$ im Untereck enthalten.

\begin{center}
\textbf{§ 201. Die Residuen.}
\end{center}

Ist $\mathfrak{P}$ ein Punkt, der $m$ mal im Untereck $\mathfrak{B}$ eines Differentials $dJ$ aufgeht ($m \mathrel{\underset{>}{=}} 0$),
so nehme man eine Funktion $z$ in $\Omega$, die in $\mathfrak{P}$ unendlich groß in der ersten Ordnung wird. Dann kann man nach § 197, (1) und § 185, (5) setzen:
\weq{1}{\frac{dJ}{dz} = a_{m-2}z^{m-2} + a_{m-3}z^{m-3} + \cdots + a_0 + a_{-1}z^{-1} + \eta z^{-2},}
worin die $a$ Konstanten sind, $\eta$ eine in $\mathfrak{P}$ endliche Funktion in $\Omega$.

\gesperrt{Der Koeffizient $-a_{-1}$ von $-z^{-1}$ heißt das Residuum des Differentials $dJ$ in bezug auf den Punkt $\mathfrak{P}$.}

Das Residuum von $dJ$ im Punkt $\mathfrak{P}$ kann nur dann von Null verschieden sein, wenn $m>0$ ist, d.\,h. wenn der Punkt $\mathfrak{P}$ im Untereck von $dJ$ vorkommt. Solche Punkte sollen Pole von $dJ$ heißen.

1. \gesperrt{Das Residuum eines eigentlichen Differentials ist gleich Null. Denn durch Differentiation einer Potenz von $z$ kann niemals $z^{-1}$ entstehen.}

Nimmt man für $z$ eine Funktion $z_1$, die gleichfalls in $\mathfrak{P}$ unendlich von der ersten Ordnung wird, so kann man setzen:
\[
z = cz_1 + \eta_1, \quad z^{-1} = c^{-1}z_1^{-1} + \eta_2,
\]
\weq{2}{\frac{dz}{dz_1} = c + z_1^{-2}\eta_3,}
worin die $c$ Konstanten, $\eta_1, \eta_2, \eta_3$, in $\mathfrak{P}$ endlich sind.

Setzt man
\[
\frac{a_{m-2}z^{m-1}}{m-1} + \frac{a_{m-3}z^{m-2}}{m-2} + \cdots + a_0 = \omega,
% [sic] printed last term $a_0$ in the expression for $\omega$, should be $a_0 z$: only then is $d\omega/dz = a_{m-2}z^{m-2} + \cdots + a_0$, as the following display for $dJ/dz_1$ requires
\]
so ergibt sich aus (1):
\[
\frac{dJ}{dz_1} = \frac{d\omega}{dz_1} + a_{-1}z^{-1}\frac{dz}{dz_1} + z^{-2}\eta',
\]
wo $\eta'$ in $\mathfrak{P}$ endlich bleibt, und darin gibt nach 1. nur der Teil
\[
a_{-1}z^{-1}\frac{dz}{dz_1}
\]
%CONT>>
